\chapter{Ce que disent les données}
\label{ch:donnees}

Les deux jeux de données ont été téléchargés et lus le 26 octobre. Les chiffres de ce chapitre sont
comptés sur les fichiers, pas recopiés d'une documentation.

\section{Banking77 : la compréhension des demandes}

\paragraph{Ce que c'est.} 13\,083 questions de clients d'une banque en ligne britannique, chacune
étiquetée par l'une de 77 intentions (Casanueva et al., 2020, licence CC-BY-4.0). Le jeu
d'entraînement compte 10\,003 questions, réparties de façon inégale (de 35 à 187 par intention) ;
le jeu de test en compte 3\,080, exactement 40 par intention. \textbf{Tout est en anglais.}

\paragraph{Les scores de référence.} L'article d'origine donne, en précision sur le test complet :
87,19\,\% pour BERT utilisé comme encodeur figé, 92,81\,\% pour USE, 93,01\,\% pour ConveRT,
93,36\,\% pour USE et ConveRT combinés, 93,66\,\% pour BERT affiné. Le message utile pour nous :
\textbf{un encodeur figé et un classifieur léger se placent à un ou deux points du meilleur modèle
affiné}, pour un coût d'entraînement de quelques minutes sur CPU.

\paragraph{Ce qu'il couvre, et surtout ce qu'il ne couvre pas.} Le vocabulaire est celui d'une
néobanque : recharges (\emph{top-up}), cartes virtuelles, Apple Pay, taux de change. Confronté aux
situations du sujet, le constat est net :

\begin{center}\small
\begin{tabularx}{\textwidth}{@{}l c Y@{}}
\toprule
\textbf{Situation du sujet} & \textbf{Couverte ?} & \textbf{Intentions Banking77 concernées} \\
\midrule
Carte perdue (Inès) & \oui & \f{lost\_or\_stolen\_card}, \f{compromised\_card},
  \f{card\_swallowed} \\
Frais (Marc) & \oui & \f{extra\_charge\_on\_statement}, \f{card\_payment\_fee\_charged},
  \f{transfer\_fee\_charged}, \f{cash\_withdrawal\_charge}\dots \\
Virement à un proche (Paul) & \partiel & Les intentions \f{transfer\_*} portent sur l'\emph{état}
  d'un virement (en attente, refusé, annulé), pas sur \og comment en faire un \fg \\
Épargne (Sophie) & \non & --- \\
Crédit auto (Karim) & \non & --- \\
Appartement (énoncé) & \non & --- \\
Prêt étudiant (oral) & \non & --- \\
Insatisfaction & \non & --- \\
Jargon (oral) & \non & --- \\
\bottomrule
\end{tabularx}
\end{center}

\begin{cle}
\textbf{Banking77 couvre le \emph{service}, pas le \emph{conseil}.} Il sait reconnaître une carte
perdue ou un frais contesté ; il ne contient pas une seule demande d'épargne, de crédit ou de projet
immobilier --- c'est-à-dire l'exemple même de l'énoncé. Les intentions de conseil sont à créer, en
français, par nous. C'est une partie du livrable \og données préparées en français \fg.
\end{cle}

\paragraph{Un bruit d'étiquetage à connaître.} Certaines étiquettes sont discutables. Dans le jeu
d'entraînement, \og \emph{What is the word?} \fg{} est rangé sous
\f{direct\_debit\_payment\_not\_recognised}, et \og \emph{Has there been any activity on my card
today?} \fg{} sous \f{lost\_or\_stolen\_card}. Cela plafonne la précision atteignable et justifie de
raisonner aussi au niveau des \emph{familles} d'intentions, plus robustes (section~\ref{sec:nlu}).

\paragraph{Neuf familles.} Nous regroupons les 77 intentions en neuf familles, vérifié par script :
chaque intention est rangée une fois et une seule. Le détail est en annexe~\ref{an:taxonomie}.

\begin{center}\small
\begin{tabular}{@{}l r r r@{}}
\toprule
\textbf{Famille} & \textbf{Intentions} & \textbf{Entraînement} & \textbf{Test} \\
\midrule
\f{carte\_perdue} & 4 & 350 & 160 \\
\f{carte\_usage} & 10 & 1\,084 & 400 \\
\f{carte\_commande} & 12 & 1\,465 & 480 \\
\f{frais\_contestation} & 14 & 2\,225 & 560 \\
\f{virement} & 10 & 1\,409 & 400 \\
\f{compte\_identite} & 8 & 884 & 320 \\
\f{especes} & 5 & 764 & 200 \\
\f{change\_etranger} & 6 & 815 & 240 \\
\f{alimenter\_compte} & 8 & 1\,007 & 320 \\
\midrule
\textbf{Total} & \textbf{77} & \textbf{10\,003} & \textbf{3\,080} \\
\bottomrule
\end{tabular}
\end{center}

\section{Bank Marketing : la base clients}

\paragraph{Ce que c'est.} Les clients d'une banque portugaise contactés par téléphone entre mai 2008
et novembre 2010 pour souscrire un dépôt à terme (UCI, Moro et al.). Le fichier
\f{bank-full.csv} compte \textbf{45\,211 clients} et 17 colonnes ; la variante
\f{bank-additional-full.csv} (41\,188 lignes) remplace le solde par des indicateurs économiques.
Nous partons de \f{bank-full.csv}, qui contient le solde.

\begin{center}\small
\begin{tabularx}{\textwidth}{@{}L{3.3cm} Y L{4.5cm}@{}}
\toprule
\textbf{Colonnes} & \textbf{Contenu} & \textbf{Usage pour nous} \\
\midrule
\f{age}, \f{job}, \f{marital}, \f{education} & Âge (18 à 95 ans, médiane 39), 12 métiers dont 938
  étudiants et 2\,264 retraités, situation (\og divorced \fg{} inclut les veufs), études &
  Profil, éligibilité (Livret Jeune, LDDS), registre de langue \\
\f{balance} & Solde annuel moyen en euros, de $-$8\,019 à 102\,127 ; médiane 448 ;
  \textbf{3\,766 soldes négatifs (8,3\,\%)} & Découverts fréquents (Marc), capacité d'épargne \\
\f{housing}, \f{loan}, \f{default} & Crédit immobilier en cours (55,6\,\%), crédit personnel
  (16,0\,\%), défaut de paiement (1,8\,\%) & Charges existantes, prudence crédit \\
\f{contact}, \f{day}, \f{month}, \f{duration}, \f{campaign}, \f{pdays}, \f{previous},
  \f{poutcome} & Historique des campagnes de démarchage & Seulement \f{previous} et
  \f{poutcome}, pour l'appétence \\
\f{y} & A souscrit un dépôt à terme : 5\,289 oui (11,7\,\%) & Cible du modèle d'appétence épargne \\
\bottomrule
\end{tabularx}
\end{center}

\begin{attention}
\textbf{La colonne \f{duration} est une fuite.} La documentation de l'UCI le dit explicitement : la
durée de l'appel n'est connue qu'\emph{après} l'appel, et elle détermine presque la réponse (un
appel de 0 seconde donne toujours \og non \fg). Un modèle qui l'utilise affiche un excellent score
et ne sert à rien. Nous l'excluons, ainsi que \f{contact}, \f{day} et \f{month}, propres à la
campagne. Le dire au jury est un point de qualité technique gratuit.
\end{attention}

\section{Ce qui manque, et comment on le comble}

Bank Marketing ne contient ni revenus, ni enfants, ni charges, ni historique d'opérations. Or
l'énoncé demande de personnaliser \og selon l'âge, les revenus et l'historique \fg. Les slides
annoncent une version \og enrichie de revenus et d'historique simulés \fg{} fournie par les
organisateurs. Deux cas :

\begin{itemize}
  \item \textbf{Elle arrive à temps.} On la branche derrière un \emph{adaptateur} unique,
        \f{src/conseiller/clients.py}, qui expose toujours le même profil au reste du code.
  \item \textbf{Elle tarde.} On simule nous-mêmes, avec une graine fixe et une méthode écrite :
        revenu mensuel net tiré par métier et âge ; nombre d'enfants à partir de l'âge et de la
        situation ; charges de crédit si \f{housing} ou \f{loan} ; douze mois d'événements (frais,
        découverts, contacts avec le service client), plus nombreux si le solde est négatif.
\end{itemize}

\begin{choix}
\textbf{Règle de transparence.} Toute donnée simulée porte le drapeau \f{simule=True} et sa méthode
est décrite dans la fiche de données. Dans l'interface, un champ simulé s'affiche en italique. Le
jury ne doit jamais pouvoir croire qu'un revenu inventé est une donnée réelle.
\end{choix}

\paragraph{Les clients de démonstration.} Les profils des slides deviennent huit clients fixes,
chacun ancré sur une ligne réelle de Bank Marketing (âge, métier, situation, solde, crédits) et
complété par simulation. Ils servent à la démo et aux tests.

\begin{center}\small
\begin{tabularx}{\textwidth}{@{}l l L{3.5cm} Y@{}}
\toprule
\textbf{Id} & \textbf{Client} & \textbf{Ancrage Bank Marketing} & \textbf{Enrichissement simulé
  (exemple)} \\
\midrule
C01 & Inès, 34 ans & \f{management}, \f{single} & 3\,200\eur{} nets, carte active \\
C02 & Marc, 52 ans & \f{blue-collar}, \f{married}, solde négatif & 1\,850\eur{} nets, 5 découverts
  en 12 mois, commissions facturées \\
C03 & Sophie, 41 ans & \f{admin.}, \f{married} & 2\,400\eur{} nets, enfants de 8 et 12 ans,
  Livret A \\
C04 & Paul, 68 ans & \f{retired}, \f{divorced} & 1\,600\eur{}, \og peu à l'aise en ligne \fg, fils
  enregistré comme bénéficiaire \\
C05 & Karim, 29 ans & \f{technician}, \f{loan=yes} & 1\,900\eur{} nets (donné), crédit en cours
  de 150\eur{}/mois \\
C06 & Julie, 37 ans & \f{services} & 3 contacts au service client aujourd'hui \\
C07 & Léa, 31 ans & \f{management}, \f{housing=no} & 2\,900\eur{} nets, 25\,000\eur{} d'épargne \\
C08 & Lucas, 20 ans & \f{student} & Job étudiant à 450\eur{}, sans épargne \\
\bottomrule
\end{tabularx}
\end{center}

\section{Le catalogue de produits}

L'énoncé le laisse à construire \og à partir d'informations publiques \fg. Nous le limitons à ce
qu'un assistant peut honnêtement proposer --- des produits \emph{simples} --- et nous l'écrivons en
JSON, une fiche par produit : identifiant, nom, catégorie, conditions d'éligibilité, plafond,
disponibilité de l'argent, à qui il convient, limites, \textbf{source publique et date de
relevé}. La liste complète et les paramètres sont en annexe~\ref{an:catalogue}.

\begin{attention}
\textbf{Les taux réglementés changent} (le taux du Livret A est révisé deux fois par an). Aucun taux
n'est écrit en dur dans le code : c'est un champ du catalogue, avec sa date et sa source, relevé le
jour même sur le site public. Les plafonds (22\,950\eur{} pour le Livret A, 12\,000\eur{} pour le
LDDS, 10\,000\eur{} pour le LEP, 1\,600\eur{} pour le Livret Jeune, 61\,200\eur{} pour le PEL) sont
stables mais sont aussi vérifiés à la source.
\end{attention}
